% ============================================
% 9. 模型灵敏度分析
% ============================================

\section{模型灵敏度分析}

为检验模型的\textbf{稳健性和可靠性}，本文对关键参数进行了灵敏度分析。通过改变\textbf{【关键参数列表】}等参数的取值，观察其对模型结果的\textbf{【影响程度】}。

\subsection{问题一的灵敏度分析}

针对问题一中的\textbf{【关键参数1】}和\textbf{【关键参数2】}，在\textbf{【原始值 $\pm 20\%$】}的范围内进行灵敏度测试，结果如图\ref{fig:sensitivity1}所示。

\begin{figure}[ht]
  \centering
  \includegraphics[width=0.8\textwidth]{../../求解/问题一/图片/【灵敏度分析图】.png}
  \caption{问题一参数灵敏度分析}
  \label{fig:sensitivity1}
\end{figure}

如图\ref{fig:sensitivity1}所示，当\textbf{【参数】}在\textbf{【范围】}内变化时，模型结果的\textbf{【变化趋势】}。具体而言，\textbf{【具体分析，≥3行】}。这表明模型对\textbf{【参数名】}的敏感程度\textbf{【高/低】}，模型的\textbf{【稳健性评价】}。

\subsection{问题二的灵敏度分析}

……

% 每个问题都必须配灵敏度分析图和 ≥3 行分析文字
